% SC26 MLIR tutorial -- Session 2, Part 1 (Linalg and tiling, CPU). Presenter deck.
%
% Build:  make            (the deck PDF)
% See presentation/README.md.
\providecommand{\notesmode}{0}
\documentclass[aspectratio=169,10pt]{beamer}
\usepackage{sc26}
\ifnum\notesmode=1 \setbeameroption{show notes}\fi

\title{Introduction to MLIR}
\subtitle{Session 2, Part 1 --- Linalg and Tiling, CPU}
\author{SC26 Tutorial}
\date{}

\begin{document}

\begin{frame}[plain]
  \titlepage
\end{frame}
\note{Pre-run setup cells 0--7 of the walkthrough notebook before the room
fills: the toolchain download is the slowest thing in this part and teaches
nothing. There is no plugin to build in this session.

One exercise, extending the tiling schedule: about 10 minutes of keyboard
time. Steps 1 and 3 are read and run together; nobody types anything in
them.}

% =======================================================================
\begin{frame}{Setup}
  \begin{itemize}
    \item<1-> \textbf{It is a fresh runtime.} Nothing from Session 1 carries
      over, so the notebook downloads the same toolchain again. Any runtime
      will do: this part runs on the CPU.
    \item<2-> \textbf{You edit one MLIR file, a schedule.} No C++ and no
      rebuild in this session: edit, save, re-run the cell.
    \item<3-> \textbf{One exercise:} tile matmul's third loop. Steps 1 and 3
      we read and run together.
    \item<4-> \textbf{Colab trick:} lost your place? \cd{Ctrl+F8} runs every
      cell above the one you are on.
  \end{itemize}
\end{frame}
\note{Matvec rides along as the second example, as in Session 1. Here it is
not an answer sheet: the schedule they extend already works on it, and the
exercise is to tile matmul's third loop.}

% =======================================================================
\begin{frame}{One level higher: Linalg}
  \begin{center}
  \begin{tikzpicture}[
      box/.style={draw=scblue,fill=scwash,rounded corners=2pt,inner sep=5pt,
                  font=\ttfamily\scriptsize,minimum height=2.4em,align=center},
      lbl/.style={font=\scriptsize\itshape,color=scgrey,align=center},
      arr/.style={-{Stealth[length=2mm]},draw=scgrey,thick}]

    \node[box,draw=scaccent] (op) at (0,0)
      {linalg.matmul\\\normalfont\itshape ships with MLIR};
    \node[box] (scf) at (5.8,0) {scf.for nest};
    \node[box] (run) at (10.8,0) {JIT + run\\\normalfont\itshape CPU, one thread};
    \draw[arr] (op) -- node[lbl,above]{convert-linalg-to-loops} (scf);
    \draw[arr] (scf) -- node[lbl,above]{Session 1's pipeline} (run);

    \visible<2->{
      \node[box,draw=scaccent] (tiled) at (2.9,-1.6)
        {linalg.matmul\\\normalfont\itshape on 32$\times$32 tiles};
      \draw[arr,draw=scaccent] (op.south) |- node[lbl,pos=0.25,left]
        {a transform\\schedule} (tiled.west);
      \draw[arr] (tiled.east) -| (scf.south);
    }
    \visible<3->{
      \node[lbl,below=0.25cm of run,font=\small\itshape,color=scaccent]
        {faster?};
    }
  \end{tikzpicture}
  \end{center}
  \vspace{0.2em}
  \begin{itemize}
    \item<1-> In Session 1 you defined \cd{blas.matmul} and wrote its
      lowering. \alert{Linalg} is a dialect that ships with MLIR: its ops
      already mean matmul and matvec, and their lowering is already written.
    \item<1-> Linalg is a more general representation of linear-algebra
      computations --- specifically \alert{einsum} computations, built on
      \alert{polyhedral} compilation methods.
    \item<2-> Linalg supports general loop transformations, like
      \alert{tiling}.
    \item<3-> Same computation, same single thread: \alert{does it get
      faster?}
  \end{itemize}
\end{frame}
\note{Step 1 is the top row: the Linalg op, lowered to loops and run, as in
Session 1. Step 2 adds the tiling. Step 3 times both and asks why.

Matmul is the theme throughout; matvec rides along as the smaller example.}

% =======================================================================
\section{Step 1 --- Linalg Ops}
\begin{frame}[plain]\sectionpage\end{frame}

\begin{frame}[fragile]{Linalg Primer: Modeling a Loop Nest}
  \begin{center}
    A bias add:\quad \cd{c[i, j] = a[i, j] + b[j]}
  \end{center}
  \vspace{0.4em}
  \begin{center}
  \begin{tikzpicture}[
      lbl/.style={font=\scriptsize,color=scink},
      mp/.style={font=\ttfamily\scriptsize,color=scaccent,fill=white,inner sep=1pt},
      arr/.style={-{Stealth[length=1.6mm]},draw=scaccent,thick}]
    % the loop nest
    \node[font=\ttfamily\footnotesize,align=left,anchor=east] at (-0.65,-0.45) {%
      \hl<1>{\lkw{for} i \lkw{in} range(4):}\\
      \hphantom{fo}\hl<1>{\lkw{for} j \lkw{in} range(8):}\\
      \hphantom{for }\hl<3>{c\hl<2>{[i, j]} = a\hl<2>{[i, j]} + b\hl<2>{[j]}}};
    % the iteration space: one point per (i, j), 4 x 8
    \foreach \i in {0,...,3} {
      \foreach \j in {0,...,7} {
        \fill[scblue!60] (\j*0.3,-\i*0.3) circle (0.045);
      }
    }
    \fill[scaccent] (1.5,-0.3) circle (0.085);
    \node[lbl] at (1.05,0.75) {the iteration space, 4 $\times$ 8};
    \draw[-{Stealth[length=1.2mm]},scgrey] (-0.05,0.25) -- node[above,lbl]{\cd{j}}
      (2.2,0.25);
    \draw[-{Stealth[length=1.2mm]},scgrey] (-0.25,0.05) -- node[left,lbl]{\cd{i}}
      (-0.25,-1.0);

    \visible<2->{
      % the data: a and c are 4 x 8, b has 8 elements; (i, j) = (1, 5)
      \fill[scaccent!45] (7.2,0.9) rectangle (7.4,1.1);
      \fill[scaccent!45] (7.2,-0.55) rectangle (7.4,-0.35);
      \fill[scaccent!45] (7.2,-1.8) rectangle (7.4,-1.6);
      \foreach \y/\h in {1.3/0.8, -0.35/0.2, -1.4/0.8} {
        \draw[step=0.2,scgrey!45,very thin,shift={(6.2,\y)}] (0,0) grid (1.6,-\h);
        \draw[scblue] (6.2,\y) rectangle (7.8,\y-\h);
      }
      \node[lbl,right] at (7.85,0.9) {\cd{a}};
      \node[lbl,right] at (7.85,-0.45) {\cd{b}};
      \node[lbl,right] at (7.85,-1.8) {\cd{c}};
      \draw[arr] (1.5,-0.3) -- node[mp,pos=0.5]{(i, j) -> (i, j)} (7.2,1.0);
      \draw[arr] (1.5,-0.3) -- node[mp,pos=0.5]{(i, j) -> (j)} (7.2,-0.45);
      \draw[arr] (1.5,-0.3) -- node[mp,pos=0.5]{(i, j) -> (i, j)} (7.2,-1.7);
    }
    \visible<3->{
      \node[lbl,color=scaccent,align=center] at (1.05,-1.65)
        {at each point:\\\texttt{c[i, j] = a[i, j] + b[j]}};
    }
  \end{tikzpicture}
  \end{center}
  \vspace{0.6em}
  \begin{itemize}
    \item<1-> The loops span an \alert{iteration space}: one point per
      \cd{(i, j)}.
    \item<2-> Each access is a \alert{map} from that space to the data: the
      point \cd{(i, j)} reads \cd{a[i, j]} and \cd{b[j]}, and writes
      \cd{c[i, j]}.
    \item<3-> Each point runs the \alert{loop body} once.
  \end{itemize}
\end{frame}
\note{No Linalg on this slide yet, just the model: the iteration space is a
box of integer points, each access is an affine map from it to an operand's
data, and the body is what one point computes. The next slide is the op that
stores exactly these three things.

The b map, (i, j) -> (j), is the broadcast: every row of the iteration space
reads the same b[j].}

% =======================================================================
\begin{frame}[fragile]{Linalg Primer: From Loops to \cd{linalg.generic}}
  \begin{columns}[t]
    \begin{column}{0.49\textwidth}
      \begin{center}\small\textbf{the loop nest}\end{center}
\begin{lstlisting}[style=half,language=mlir]
(@\hl<2>{\lkw{scf.for} \%i = \%c0 \lkw{to} \%c4 \lkw{step} \%c1 \{}@)
  (@\hl<2>{\lkw{scf.for} \%j = \%c0 \lkw{to} \%c8 \lkw{step} \%c1 \{}@)
    %x = memref.load %a(@\hl<3>{[\%i, \%j]}@) : memref<4x8xf32>
    %y = memref.load %b(@\hl<3>{[\%j]}@) : memref<8xf32>
    (@\hl<4>{\%s = \lkw{arith.addf} \%x, \%y : \lkw{f32}}@)
    memref.store %s, %c(@\hl<3>{[\%i, \%j]}@) : memref<4x8xf32>
  }
}
\end{lstlisting}
    \end{column}
    \begin{column}{0.49\textwidth}
      \begin{center}\small\textbf{the same, as \cd{linalg.generic}}\end{center}
\begin{lstlisting}[style=half,language=mlir]
(@\hl<3>{\#id = \lkw{affine\_map}<(i, j) -> (i, j)>}@)
(@\hl<3>{\#bcast = \lkw{affine\_map}<(i, j) -> (j)>}@)
linalg.generic {
    (@\hl<3>{indexing\_maps = [\#id, \#bcast, \#id]}@),
    (@\hl<2>{iterator\_types = ["parallel", "parallel"]}@)}
    (@\hl<5>{ins(\%a, \%b : \lkw{memref}<4x8xf32>, \lkw{memref}<8xf32>)}@)
    (@\hl<5>{outs(\%c : \lkw{memref}<4x8xf32>)}@) {
(@\hl<4>{\textasciicircum{}bb0(\%x: \lkw{f32}, \%y: \lkw{f32}, \hl<5>{\%out: \lkw{f32}}):}@)
  (@\hl<4>{\%s = \lkw{arith.addf} \%x, \%y : \lkw{f32}}@)
  (@\hl<4>{\lkw{linalg.yield} \%s : \lkw{f32}}@)
}
\end{lstlisting}
    \end{column}
  \end{columns}
  \vspace{0.5em}
  \begin{itemize}
    \item<2-> \cd{iterator\_types}: one per loop, and the dependence between
      points. Both are \cd{parallel}: no point depends on another.
    \item<3-> \cd{indexing\_maps}: one per operand, the maps from the last
      slide.
    \item<4-> The region is the loop body, run at each point: it receives the
      elements and yields the result.
    \item<5-> \cd{ins} are read; \cd{outs} is written. The region also
      receives the output's current value, \cd{\%out}, so an op can
      accumulate into it.
  \end{itemize}
\end{frame}
\note{Both sides are valid MLIR, and convert-linalg-to-loops turns the right
one into the left one (with generated names). affine\_map dimensions can be
named anything; mlir-opt prints them back as d0, d1.

The generic op says nothing about loop order or how the loops are run; it
only says what each point of the iteration space computes. That is what
leaves the compiler free to tile it.

The bias add ignores \%out, which is why its loop nest never loads c. Matmul
and matvec use it: that is their +=.}

% =======================================================================
\begin{frame}[fragile]{Named Linalg Ops}
  In \cd{linalg\_matmul.mlir} and \cd{linalg\_matvec.mlir}:
\begin{lstlisting}[language=mlir]
linalg.matmul ins(%a, %b : ...) outs(%c : ...)
linalg.matvec ins(%mat, %vec : ...) outs(%out : ...)
\end{lstlisting}
  \vspace{0.2em}
  \begin{itemize}
    \item<2-> A \alert{named} Linalg op is really a \cd{linalg.generic} whose
      indexing maps, iterator types and region are fixed by its name:
  \end{itemize}
  \vspace{0.2em}
  \visible<2->{
  \begin{center}\scriptsize
    \renewcommand{\arraystretch}{1.2}
    \begin{tabular}{lll}
      & \texttt{linalg.matmul} & \texttt{linalg.matvec} \\
      \hline
      \texttt{indexing\_maps} & \texttt{a:\ (i, j, k) -> (i, k)}
        & \texttt{mat:\ (i, j) -> (i, j)} \\
      & \texttt{b:\ (i, j, k) -> (k, j)} & \texttt{vec:\ (i, j) -> (j)} \\
      & \texttt{c:\ (i, j, k) -> (i, j)} & \texttt{out:\ (i, j) -> (i)} \\
      \hline
      \texttt{iterator\_types} & \texttt{parallel, parallel, reduction}
        & \texttt{parallel, reduction} \\
      \hline
      region yields & \texttt{c + a * b} & \texttt{out + mat * vec} \\
    \end{tabular}
  \end{center}}
  \vspace{0.2em}
  \begin{itemize}
    \item<3-> Running \cd{linalg-generalize-named-ops} rewrites each named op
      as the \cd{linalg.generic} it stands for: that is where this table
      comes from.
    \item<4-> All Linalg named ops and \cd{linalg.generic} implement the
      \alert{\cd{LinalgOp} interface}, and Linalg transformations are written
      against that interface, so any transformation applies to all of them.
  \end{itemize}
\end{frame}
\note{Both files are complete; there is nothing to write in step 1. Round-trip
both in the notebook (cell 11) as you show them. The listing leaves out the
operand types: they are the 128x512, 512x256 and 128x256 memrefs (512 and 128
for matvec's vector and output).

Run cell 13 on the third step: it is linalg-generalize-named-ops on both
files. mlir-opt names the dimensions d0, d1, d2 rather than i, j, k.

The LinalgOp interface is what step 2's schedule matches on
(structured.match interface\{LinalgOp\}), which is why it never has to name
an op.

reduction: points along k (along j for matvec) add into the same output
element. The next slide shows what that does to the loops. Do not explain it
as the output's map dropping k.

Every Linalg named op on memrefs accumulates into its outs operand. That is
why run\_and\_check zeroes the output buffer before each call.}

% =======================================================================
\begin{frame}[fragile]{Lowered to Loops}
  \cd{mlir-opt linalg\_matvec.mlir -{}-convert-linalg-to-loops}, the body of
  \cd{@matvec}:
\begin{lstlisting}[language=mlir]
%c0 = arith.constant 0 : index
%c128 = arith.constant 128 : index
%c1 = arith.constant 1 : index
%c512 = arith.constant 512 : index
(@\hl<2>{\lkw{scf.for} \%arg3 = \%c0 \lkw{to} \%c128 \lkw{step} \%c1 \{}@)                      // i
  (@\hl<2>{\lkw{scf.for} \%arg4 = \%c0 \lkw{to} \%c512 \lkw{step} \%c1 \{}@)                    // j
    (@\hl<3>{\%0 = \lkw{memref.load} \%arg0[\%arg3, \%arg4] : \lkw{memref}<128x512xf32>}@)
    (@\hl<3>{\%1 = \lkw{memref.load} \%arg1[\%arg4] : \lkw{memref}<512xf32>}@)
    (@\hl<4>{\%2 = \lkw{memref.load} \%arg2[\%arg3] : \lkw{memref}<128xf32>}@)
    (@\hl<5>{\%3 = \lkw{arith.mulf} \%0, \%1 : \lkw{f32}}@)
    (@\hl<5>{\%4 = \lkw{arith.addf} \%2, \%3 : \lkw{f32}}@)
    (@\hl<5>{\lkw{memref.store} \%4, \%arg2[\%arg3] : \lkw{memref}<128xf32>}@)
  }
}
\end{lstlisting}
  \vspace{0.2em}
  \begin{itemize}
    \item<2-> One \cd{scf.for} per iterator type: \cd{i}, \cd{j}.
    \item<3-> A load of each input, at the indices its map selects.
    \item<4-> A load of \cd{out} too: \cd{j} is a \alert{reduction}, so every
      point along \cd{j} adds into the same \cd{out[i]}, and the body reads
      its current value first.
    \item<5-> The body, then a store back into \cd{out}. It is Session 1's
      \cd{matvec\_scf.mlir}, generated.
  \end{itemize}
\end{frame}
\note{The // i, j comments are the only additions to the real output.

The bias add from the primer, both loops parallel, lowers with no load of c
at all. Strictly, the load is there because the body uses the output's
current value, which a reduction needs. Do not explain it as out's map
dropping j.

One difference from Session 1's matvec\_scf.mlir, if someone spots it: that
file loads out[i] once and carries the sum in iter\_args; this nest loads and
stores out[i] on every j iteration. Same answer; LLVM cleans it up.}

% =======================================================================
\begin{frame}[fragile]{Lowering and Execution}
\begin{lstlisting}[style=plain,basicstyle=\ttfamily\small]
(@\hl<2>{convert-linalg-to-loops}@)
(@{\normalfont\itshape then Session 1's CPU lowering pipeline}@)
\end{lstlisting}
  \vspace{0.3em}
  \begin{itemize}
    \item<2-> \cd{convert-linalg-to-loops} does the job your
      \cd{convert-blas-to-scf} did in Session 1.
    \item<3-> JIT-compiled, and checked against numpy's \cd{a @ b}.
  \end{itemize}
  \vspace{0.2em}
  \visible<3->{\begin{center}\color{scgreen}\large\ttfamily OK\end{center}}
\end{frame}
\note{Cell 17. Its CPU\_LOWERING\_PIPELINE is llvm-request-c-wrappers,
convert-linalg-to-loops, then the rest of Session 1's passes:
convert-scf-to-cf, convert-cf-to-llvm, convert-arith-to-llvm,
finalize-memref-to-llvm, convert-func-to-llvm, reconcile-unrealized-casts.
In Session 1 convert-blas-to-scf ran as its own step before the pipeline;
here the Linalg lowering is simply part of it.

The run is @matmul, as in the notebook. Step 2 adds two passes to this
pipeline for the tiled matmul.}

% =======================================================================
\section{Step 2 --- Tiling, with a Transform Schedule}
\begin{frame}[plain]\sectionpage\end{frame}

\begin{frame}{Tiling}
  \begin{center}
  \begin{tikzpicture}[
      lbl/.style={font=\scriptsize,color=scink},
      sym/.style={font=\large,color=scgrey}]
    % c += a x b to scale, one grid cell per 32 x 32 tile:
    % c is 128 x 256 (4 x 8 tiles), a 128 x 512 (4 x 16), b 512 x 256 (16 x 8)
    \fill[scaccent!70] (1.0,0) rectangle (1.2,0.2);
    \visible<2->{
      \fill[scaccent!35] (2.6,0) rectangle (5.8,0.2);
      \fill[scaccent!35] (7.8,-1.6) rectangle (8.0,1.6);
    }
    \visible<3->{
      % one k step: the fourth 32 x 32 tile of each strip
      \fill[scaccent!70] (3.2,0) rectangle (3.4,0.2);
      \fill[scaccent!70] (7.8,0.8) rectangle (8.0,1.0);
    }
    \foreach \x/\w/\h in {0/1.6/0.4, 2.6/3.2/0.4, 6.8/1.6/1.6} {
      \draw[step=0.2,scgrey!45,very thin] (\x,-\h) grid (\x+\w,\h);
      \draw[scblue] (\x,-\h) rectangle (\x+\w,\h);
    }
    \visible<3->{
      \draw[-{Stealth[length=1.2mm]},scaccent] (2.6,0.6) --
        node[above,lbl,color=scaccent]{\texttt{k}} (5.8,0.6);
      \draw[-{Stealth[length=1.2mm]},scaccent] (8.6,1.6) --
        node[right,lbl,color=scaccent]{\texttt{k}} (8.6,-1.6);
    }
    \node[lbl] at (0.8,-0.65) {\cd{c}, 128 $\times$ 256};
    \node[lbl] at (4.2,-0.65) {\cd{a}, 128 $\times$ 512};
    \node[lbl] at (7.6,-1.85) {\cd{b}, 512 $\times$ 256};
    \node[sym] at (2.1,0) {\texttt{+=}};
    \node[sym] at (6.3,0) {$\times$};
    \draw[-{Stealth[length=1.5mm]},scaccent] (1.1,1.25) node[above,lbl]
      {one 32 $\times$ 32 tile of \cd{c}} -- (1.1,0.25);
  \end{tikzpicture}
  \end{center}
  \vspace{0.1em}
  \begin{itemize}
    \item<1-> Tiling splits a loop nest into two sets of loops:
      \begin{itemize}
        \item Outer loops that step from one tile to the next;
        \item Inner loops that do the computation within a tile;
        \item The same computation, in a different order.
      \end{itemize}
    \item<2-> Tile \cd{i} and \cd{j} by 32, and each 32 $\times$ 32 tile of
      \cd{c} reads a 32-row strip of \cd{a} and a 32-column strip of \cd{b}.
    \item<3-> Tile \cd{k} by 32 too, and each strip is read one
      32 $\times$ 32 tile at a time: 16 steps along \cd{k} for each tile of
      \cd{c}.
  \end{itemize}
\end{frame}
\note{Exactly the same, to the last bit: tiling i and j does not reorder any
sum, and tiling k keeps the additions for each element in k order. That is
why the notebook checks the tiled result against numpy with the same
tolerance as the untiled one.

The order here, i and j first and then k inside the strips, is the order the
exercise's schedule tiles in. The solution's subview of a subview is this
picture in IR.}

% =======================================================================
\begin{frame}[fragile]{The Tile Is Still a \cd{linalg.matmul}}
  \begin{columns}[t]
    \begin{column}{0.35\textwidth}
      \begin{center}\small\textbf{the op we started with}\end{center}
\begin{lstlisting}[style=half,language=mlir]
linalg.matmul
    ins(%a, %b : memref<128x512xf32>,
                 memref<512x256xf32>)
    outs(%c : memref<128x256xf32>)
\end{lstlisting}
    \end{column}
    \begin{column}{0.63\textwidth}
      \begin{center}\small\textbf{tiled: the shape of it}\end{center}
\begin{lstlisting}[style=half,language=mlir]
(@\hl<2>{\lkw{scf.for} \%i = 0 \lkw{to} 128 \lkw{step} 32 \{}@)              // tile i
  (@\hl<2>{\lkw{scf.for} \%j = 0 \lkw{to} 256 \lkw{step} 32 \{}@)            // tile j
    (@\hl<2>{\lkw{scf.for} \%k = 0 \lkw{to} 512 \lkw{step} 32 \{}@)          // tile k
      (@\hl<3>{\%a\_tile = \lkw{memref.subview} \%a[\%i, \%k] [32, 32]}@)
      (@\hl<3>{\%b\_tile = \lkw{memref.subview} \%b[\%k, \%j] [32, 32]}@)
      (@\hl<3>{\%c\_tile = \lkw{memref.subview} \%c[\%i, \%j] [32, 32]}@)
      (@\hl<4>{\lkw{linalg.matmul} ins(\%a\_tile, \%b\_tile) outs(\%c\_tile)}@)
    }
  }
}
\end{lstlisting}
    \end{column}
  \end{columns}
  \vspace{0.4em}
  \begin{itemize}
    \item<2-> The outer loops are \cd{scf.for}s over \cd{i}, \cd{j} and
      \cd{k}, each stepping by 32.
    \item<3-> \cd{memref.subview} picks out a 32 $\times$ 32 tile of each
      operand.
    \item<4-> The inner loops are still one \cd{linalg.matmul}: the op we
      started with, on 32 $\times$ 32 tiles.
  \end{itemize}
\end{frame}
\note{The shape, not real output: the names are made up, and constants and
types are left out. The real IR is nested differently: the exercise's
schedule tiles i and j first and k second, so its k loop slices the strips
the first step made rather than a, b and c directly. The solution, a few
slides on, shows it.

The point to land: tiling does not lower anything. What sits inside the
tile loops is the same Linalg op, so it can be tiled again or lowered the
usual way.}

% =======================================================================
\begin{frame}[fragile]{A Schedule: MLIR That Transforms MLIR}
  \begin{center}
  \begin{tikzpicture}[
      box/.style={draw=scblue,fill=scwash,rounded corners=2pt,inner sep=5pt,
                  font=\ttfamily\scriptsize,minimum height=2.4em,align=center},
      lbl/.style={font=\scriptsize\itshape,color=scgrey,align=center},
      arr/.style={-{Stealth[length=2mm]},draw=scgrey,thick}]
    \node[box,draw=scaccent] (sch) at (0,0.65)
      {tile\_linalg.mlir\\\normalfont\itshape the schedule};
    \node[box] (pay) at (0,-0.65)
      {linalg\_matmul.mlir\\\normalfont\itshape the payload};
    \node[box] (int) at (5.0,0)
      {mlir-opt\\\normalfont\itshape transform interpreter};
    \node[box] (out) at (9.8,0)
      {the payload,\\\normalfont\itshape tiled};
    \draw[arr] (sch.east) -- (int.west);
    \draw[arr] (pay.east) -- (int.west);
    \draw[arr] (int) -- (out);
  \end{tikzpicture}
  \end{center}
  \vspace{0.2em}
  \begin{itemize}
    \item<1-> In Session 1 a rewrite meant C++: a pattern, a plugin build,
      and reloading \cd{mlir-opt}.
    \item<2-> The \alert{transform dialect} lets you specify transformations
      on operations, and apply them.
    \item<3-> Its operations are transformations to apply. Its values are
      \alert{handles} to operations, and an operation's results are handles
      to the operations it creates.
    \item<4-> No C++, no rebuild: edit the schedule, re-run the cell.
  \end{itemize}
\end{frame}
\note{The notebook runs it as
mlir-opt linalg\_matmul.mlir
-{}-transform-preload-library=transform-library-paths=tile\_linalg.mlir
-{}-transform-interpreter. The preload flag loads the schedule; the
interpreter finds @\_\_transform\_main in it and applies it to the payload.

Keep the schedule in its own file. A payload module that carries the
transform attribute trips the transform verifier as soon as a later pass
outlines a function into it.}

% =======================================================================
\begin{frame}[fragile]{Reference: \cd{tile\_linalg.mlir}, Step 1 of 3 --- Schedule}
\begin{lstlisting}[language=mlir]
module attributes {transform.with_named_sequence} {
  (@\hl<2>{\lkw{transform.named\_sequence} @\_\_transform\_main(}@)
      (@\hl<2>{\%root: !transform.any\_op \{transform.readonly\}) \{}@)
    (@\hl<3>{\%ops = \lkw{transform.structured.match} interface\{LinalgOp\} in \%root}@)
      : (!transform.any_op) -> !transform.any_op
    (@\hl<5>{\%tiled, \%loops:2}@) = (@\hl<4>{\lkw{transform.structured.tile\_using\_for} \%ops}@)
      (@\hl<4>{tile\_sizes [32, 32]}@)
      : (!transform.any_op) -> ((@\hl<5>{!transform.any\_op, !transform.any\_op,}@)
                                (@\hl<5>{!transform.any\_op}@))
    transform.yield
  }
}
\end{lstlisting}
  \vspace{0.3em}
  \begin{itemize}
    \item<2-> The interpreter calls \cd{@\_\_transform\_main}, passing the
      payload's top-level op as \cd{\%root}.
    \item<3-> \cd{match} finds ops in the payload --- by \alert{interface},
      not by name, so any Linalg op. \cd{\%ops} is a \alert{handle}: a
      reference to payload ops, not a value in your program.
    \item<4-> \cd{tile\_using\_for} takes one tile size per loop, in
      \cd{iterator\_types} order. \cd{[32, 32]} tiles the first two.
    \item<5-> It returns the tiled op, then one handle per loop it created:
      \cd{\%loops:2}.
  \end{itemize}
\end{frame}
\note{session2/tile\_linalg.mlir, as in the repo without its comments. It is a
working schedule; they read it before changing anything.

tile\_using\_for is itself written against the LinalgOp interface, which is
why one schedule tiles matvec and matmul alike.}

% =======================================================================
\begin{frame}[fragile]{Reference: \cd{tile\_linalg.mlir}, Step 2 of 3 --- Applied to Matvec}
\begin{lstlisting}[language=mlir,basicstyle=\ttfamily\fontsize{6.6}{7.9}\selectfont]
(@\hl<2>{\lkw{scf.for} \hl<4>{\%arg3} = \%c0 \lkw{to} \%c128 \lkw{step} \%c32 \{}@)                                 // i, by 32
  (@\hl<2>{\lkw{scf.for} \hl<4>{\%arg4} = \%c0\_0 \lkw{to} \%c512 \lkw{step} \%c32\_1 \{}@)                          // j, by 32
    %subview = memref.subview %arg0(@\hl<4>{[\%arg3, \%arg4]}@) (@\hl<3>{[32, 32]}@) [1, 1] : ... to memref<32x32xf32, ...>
    %subview_2 = memref.subview %arg1(@\hl<4>{[\%arg4]}@) (@\hl<3>{[32]}@) [1] : ... to memref<32xf32, ...>
    %subview_3 = memref.subview %arg2(@\hl<4>{[\%arg3]}@) (@\hl<3>{[32]}@) [1] : ... to memref<32xf32, ...>
    (@\hl<5>{\lkw{linalg.matvec} ins(\%subview, \%subview\_2 : ...) outs(\%subview\_3 : ...)}@)
  }
}
\end{lstlisting}
  \vspace{0.4em}
  \begin{itemize}
    \item<2-> Two \cd{scf.for} tile loops stepping by 32: both of matvec's
      loops are tiled.
    \item<3-> \cd{memref.subview} takes a 32-wide \alert{view} of each
      operand: offsets, sizes, strides. Nothing is copied.
    \item<4-> Each slice follows its operand's indexing map. \cd{vec}'s is
      \cd{(i, j) -> (j)}, so its slice moves with the \cd{j} loop
      (\cd{\%arg4}); \cd{out}'s is \cd{(i, j) -> (i)}, so it moves with the
      \cd{i} loop (\cd{\%arg3}). \cd{mat}'s is the identity: it moves with
      both.
    \item<5-> Inside, the original op on those slices. Still Linalg: tiling
      restructures the loops, it lowers nothing.
  \end{itemize}
\end{frame}
\note{The real output of the notebook's cell 23, compressed: the constants
are left out, and each "..." is a type the slide has no room for. The slice
types carry a layout, strided<[512, 1], offset: ?>: a slice keeps its
parent's row stride.}

% =======================================================================
\begin{frame}[fragile]{Reference: \cd{tile\_linalg.mlir}, Step 3 of 3 --- Applied to Matmul}
\begin{lstlisting}[language=mlir,basicstyle=\ttfamily\fontsize{6.6}{7.9}\selectfont]
scf.for %arg3 = %c0 to %c128 step %c32 {                                 // i, by 32
  scf.for %arg4 = %c0_0 to %c256 step %c32_1 {                           // j, by 32
    %subview = memref.subview %arg0[%arg3, 0] [32, (@\hl<2>{512}@)] [1, 1] : ... to memref<32x512xf32, ...>
    %subview_2 = memref.subview %arg1[0, %arg4] [(@\hl<2>{512}@), 32] [1, 1] : ... to memref<512x32xf32, ...>
    %subview_3 = memref.subview %arg2[%arg3, %arg4] [32, 32] [1, 1] : ... to memref<32x32xf32, ...>
    linalg.matmul ins(%subview, %subview_2 : ...) outs(%subview_3 : ...)
  }
}
\end{lstlisting}
  \vspace{0.4em}
  \begin{itemize}
    \item<1-> The same schedule, unchanged: \cd{i} and \cd{j} tiled by 32,
      as for matvec.
    \item<2-> Matmul has a third loop, \cd{k}, and \cd{[32, 32]} leaves it
      alone. The slices of \cd{a} and \cd{b} are 32 $\times$ \alert{512} and
      \alert{512} $\times$ 32: the 512 is \cd{k}, still untiled.
  \end{itemize}
\end{frame}
\note{The real output of cell 25, compressed the same way. These are the
strips from the tiling slide: a 32-row strip of a, a 32-column strip of b, a
32 by 32 tile of c.}

% =======================================================================
\begin{frame}[fragile]{Exercise 6 --- Tile Matmul's \cd{k} Loop}
  \begin{center}\large
    Extend \cd{session2/tile\_linalg.mlir} so that matmul's \cd{k} loop is
    tiled by 32 too.
  \end{center}
  \vspace{0.1em}
  \textbf{What is different from matvec:}
  \begin{itemize}
    \item Matmul has three loops; matvec has two.
    \item Two things you may need: \cd{transform.structured.match
      ops\{["linalg.matmul"]\}} matches by op name, and a tile size of
      \cd{0} leaves that loop untiled.
  \end{itemize}
  \vspace{0.4em}
  \textbf{Check as you go --- round-trip it} after every save:
\begin{lstlisting}[style=plain,basicstyle=\ttfamily\small]
!mlir-opt tile_linalg.mlir
\end{lstlisting}
  \vspace{0.1em}
  \begin{itemize}
    \item That checks the schedule is well-formed, not that it does what you
      meant. For that, re-run the matmul cell to apply it, then the check
      after it: matmul's innermost op must work on 32 $\times$ 32 slices.
  \end{itemize}
  \vspace{0.4em}
  \begin{center}\color{scaccent}\large
    10 minutes to write it --- then we walk through the solution.
  \end{center}
\end{frame}
\note{Changing the shared step to [32, 32, 32] needs \%loops:3 as well. With
\%loops:2 left as it was, the round-trip rejects it: "expected number of
loops to tile (3) to match number of `loops` results (2)". With \%loops:3 it
round-trips and tiles matmul, but the notebook's matvec cell (23), which runs
after the exercise, fails on it: "too many tiles provided, expected at most 2
found 3".

After an edit, send them back to the matmul cell (25) and then the check
(27). Not Ctrl+F8 on the check: that re-runs everything above it, the matvec
cell included.

Until the exercise is done, the notebook stops at the k-loop check (cell 27):
"@matmul is not tiled in all three loops yet".}

% =======================================================================
\begin{frame}[fragile]{Solution --- \cd{tile\_linalg.mlir}, and What Changed}
\begin{lstlisting}[language=mlir,basicstyle=\ttfamily\fontsize{6.6}{7.9}\selectfont]
module attributes {transform.with_named_sequence} {
  transform.named_sequence @__transform_main(
      %root: !transform.any_op {transform.readonly}) {
    %ops = transform.structured.match interface{LinalgOp} in %root
      : (!transform.any_op) -> !transform.any_op
    %tiled, %loops:2 = transform.structured.tile_using_for %ops
      tile_sizes [32, 32]
      : (!transform.any_op) -> (!transform.any_op, !transform.any_op,
                                !transform.any_op)
    (@\df{\%matmul = transform.structured.match ops\{["linalg.matmul"]\} in \%root}@)
      (@\df{: (!transform.any\_op) -> !transform.any\_op}@)
    (@\df{\%tiled\_k, \%k\_loop = transform.structured.tile\_using\_for \%matmul}@)
      (@\df{tile\_sizes [0, 0, 32]}@)
      (@\df{: (!transform.any\_op) -> (!transform.any\_op, !transform.any\_op)}@)
    transform.yield
  }
}
\end{lstlisting}
  \vspace{0.2em}
  \begin{itemize}
    \item<2-> You could also try \cd{[32, 32, 32]} in the first step.
  \end{itemize}
\end{frame}
\note{The second match runs after the first tiling, so it finds the new
linalg.matmul inside the i and j tile loops, and tiles that.

If someone does try [32, 32, 32] in the first step: it tiles matmul, but
the notebook also applies the schedule to matvec, which stops with "too many
tiles provided, expected at most 2 found 3". The exercise slide's note has
both errors.}

% =======================================================================
\begin{frame}[fragile]{Solution --- Tiled Matmul}
\begin{lstlisting}[language=mlir,basicstyle=\ttfamily\fontsize{6.4}{7.6}\selectfont]
scf.for %arg3 = %c0 to %c128 step %c32 {                                  // i, by 32
  scf.for %arg4 = %c0_0 to %c256 step %c32_1 {                            // j, by 32
    (@\hl<3>{\%subview}@) = memref.subview %arg0[%arg3, 0] [32, 512] [1, 1] : ...
    (@\hl<3>{\%subview\_2}@) = memref.subview %arg1[0, %arg4] [512, 32] [1, 1] : ...
    %subview_3 = memref.subview %arg2[%arg3, %arg4] [32, 32] [1, 1] : ...
    ...
    (@\hl<2>{\lkw{scf.for} \%arg5 = \%c0\_4 \lkw{to} \%c512 \lkw{step} \%c32\_5 \{}@)                       // k, by 32
      %subview_6 = memref.subview (@\hl<3>{\%subview}@)[0, %arg5] (@\hl<2>{[32, 32]}@) [1, 1] : ...
      %subview_7 = memref.subview (@\hl<3>{\%subview\_2}@)[%arg5, 0] (@\hl<2>{[32, 32]}@) [1, 1] : ...
      %subview_8 = memref.subview %subview_3[0, 0] [32, 32] [1, 1] : ...
      (@\hl<4>{\lkw{linalg.matmul} ins(\%subview\_6, \%subview\_7 : \lkw{memref}<32x32xf32, ...>, \lkw{memref}<32x32xf32, ...>)}@)
                    (@\hl<4>{outs(\%subview\_8 : \lkw{memref}<32x32xf32, ...>)}@)
    }
  }
}
\end{lstlisting}
  \vspace{0.2em}
  \begin{itemize}
    \item<2-> A third tile loop, over \cd{k}, stepping by 32.
    \item<3-> It takes a \alert{subview of a subview}: \cd{\%subview\_6} is cut
      from \cd{\%subview}, the 32 $\times$ 512 strip of \cd{a} that tiling
      \cd{i} and \cd{j} made, not from \cd{a} itself. The same for \cd{b}.
    \item<4-> The check: every operand of the innermost \cd{linalg.matmul} is
      32 $\times$ 32.
  \end{itemize}
\end{frame}
\note{The real output with the solved schedule, compressed: the "..." after
subview\_3 is three constants, the others are types. Run cells 25 to 29 in
the walkthrough notebook as you go.

Only the innermost op tells you all three loops are tiled. The 32 by 512
subviews from the first step are still in the IR when the exercise is done,
which is why the check looks at the linalg.matmul line and nothing else.}

% =======================================================================
\begin{frame}[fragile]{Lowering the Tiled Matmul}
\begin{lstlisting}[style=plain,basicstyle=\ttfamily\scriptsize]
convert-linalg-to-loops
(@\df{expand-strided-metadata}@)    (@\only<2>{\color{scaccent}\frenchspacing $\leftarrow$ subviews become pointer and stride arithmetic \dots}@)
(@\df{lower-affine}@)               (@\only<3>{\color{scaccent}\frenchspacing $\leftarrow$ \dots{} which comes out as affine.apply: make it arith}@)
(@{\normalfont\itshape then Session 1's CPU lowering pipeline}@)
\end{lstlisting}
  \vspace{0.3em}
  \begin{itemize}
    \item<1-> Tiling introduced \cd{memref.subview}, which the step 1 pipeline
      never saw. Two more passes lower it, right after
      \cd{convert-linalg-to-loops}.
    \item<4-> Same numpy check as step 1:
  \end{itemize}
  \vspace{0.2em}
  \visible<4->{\begin{center}\color{scgreen}\large\ttfamily OK\end{center}}
\end{frame}
\note{Cell 29 redefines CPU\_LOWERING\_PIPELINE with the two passes, then runs
the tiled matmul. The later cells that compile a file use it too, which does
nothing to the untiled one.

Leave expand-strided-metadata out and the JIT stops on a bare "LLVM
Translation failed for operation: builtin.unrealized\_conversion\_cast".
Leave lower-affine out after it and the ExecutionEngine fails to build with
no further detail.}

% =======================================================================
\section{Step 3 --- Did Matmul Get Faster?}
\begin{frame}[plain]\sectionpage\end{frame}

\begin{frame}{Timing It}
  \begin{itemize}
    \item<1-> Same computation, same numbers, same single thread: untiled
      against tiled.
    \item<2-> \cd{benchmark} keeps the \alert{fastest} of five runs. Colab is
      a shared machine, so a slow run mostly measures your neighbours.
  \end{itemize}
  \vspace{1em}
  \visible<3->{
  \begin{center}\large
    \begin{tabular}{lr}
      untiled & 9.0 ms \\
      tiled   & 0.8 ms \\
      \hline
      speedup & \alert{10.9$\times$} \\
    \end{tabular}
  \end{center}
  \begin{center}\scriptsize\color{scgrey}
    one run on a desktop CPU --- yours will differ
  \end{center}}
\end{frame}
\note{Run cell 31 live and read out the room's numbers rather than these.
These are from the presenter's machine (a Threadripper 7960X); the shapes
changed to 128, 256, 512 after the last Colab measurement, so Colab's figure
for them is not known yet.}

% =======================================================================
\begin{frame}{Why So Much Faster? Read the Machine Code}
  \begin{itemize}
    \item<1-> The usual answer is the cache. Not here: the three matrices
      together are under 1 MB, small enough to sit largely in cache untiled.
    \item<2-> What changed is the \alert{code LLVM generated}. The next cell
      dumps the JIT's object code and counts the float multiplies and adds in
      \cd{matmul}:
  \end{itemize}
  \vspace{0.3em}
  \visible<2->{
  \begin{center}
    \begin{tabular}{lll}
      \textbf{scalar} & \cd{mulss}, \cd{addss}, \dots & one float at a time \\
      \textbf{packed} & \cd{mulps}, \cd{addps}, \dots &
        4, 8 or 16 floats at once, in \cd{xmm}, \cd{ymm}, \cd{zmm} \\
    \end{tabular}
  \end{center}}
  \vspace{0.3em}
  \begin{itemize}
    \item<3-> Untiled: \alert{scalar only}. Tiled: \alert{packed}, 8 or 16
      floats at once, depending on the CPU.
  \end{itemize}
\end{frame}
\note{Cell 33 installs objdump if it is missing; cell 34 is the count. On the
presenter's machine (AVX-512) the tiled matmul reported 16 floats at once;
the Colab CPUs report 8 or 16.

Why packed, if someone asks: in a tile, k runs a constant 32 times and LLVM
unrolls it completely. The j loop is then innermost, its iterations are
independent, and LLVM vectorizes it. Untiled, the innermost loop is the k
reduction, a running float sum; vectorizing it would reorder the additions,
so LLVM does not. Both are cost-model decisions: tiles of 16 and 32 get
through on every CPU Colab hands out, 24 and 64 stay scalar
(session2/README.md has the table).}

% =======================================================================
\begin{frame}{Recap}
  \begin{center}
  \begin{tikzpicture}[
      box/.style={draw=scblue,fill=scwash,rounded corners=2pt,inner sep=5pt,
                  font=\ttfamily\scriptsize,minimum height=2.4em,align=center},
      lbl/.style={font=\scriptsize\itshape,color=scgrey,align=center},
      arr/.style={-{Stealth[length=2mm]},draw=scgrey,thick}]
    \node[box] (op) at (0,0) {linalg.matmul};
    \node[box,draw=scaccent] (tiled) at (4.6,0)
      {linalg.matmul\\\normalfont\itshape on 32$\times$32 tiles};
    \node[box] (run) at (10.0,0)
      {JIT + run\\\normalfont\itshape CPU, one thread};
    \draw[arr,draw=scaccent] (op) -- node[lbl,above]{tile\_linalg.mlir} (tiled);
    \draw[arr] (tiled) -- node[lbl,above]{loops + Session 1's\\pipeline} (run);
  \end{tikzpicture}
  \end{center}
  \vspace{0.3em}
  \begin{itemize}
    \item<2-> Linalg ops describe a computation --- loops, maps, body ---
      through one interface, \cd{LinalgOp}.
    \item<3-> A schedule written against that interface tiles any Linalg op.
      No C++, no rebuild.
  \end{itemize}
\end{frame}
\note{End here.}

\end{document}
